%=============================================================================!
% 17.3  A LEARNED POTENTIAL IN BOTH                                           !
%=============================================================================!
\section{A learned potential in both}
\label{sec:ip-learned}

A translator who has read millions of sentences can render one never
seen before, but only in the languages that reading covered, and the
gaps in a translation show where the reading was thin. A foundation
model of interatomic forces, one trained once on a wide range of
materials to be used on any of them, is in the same position.
MACE-MP-0\cite{Batatia2025} is a network of the kind that the planned continuation of Part~V builds,
trained on the calculations of the Materials Project\cite{Deng2023},
compounds of 89 elements, at every structure each calculation passed
through on its way to its least energy. They use density-functional
theory (\cref{sec:pe-origin}), which needs an approximation, called the
functional, for one part of the electrons' energy; theirs is
PBE\cite{Perdew1996}, with an added correction for the oxides and
fluorides of some transition metals. The model gives energies and forces
for structures it was never shown. This section runs it in ASE
and in LAMMPS, measures what it costs, finds one of its gaps, and chooses
the time step for lithium in graphite.

\subsection*{Running it}

In ASE the model is a calculator like any other: MACE's package provides
one that loads the trained weights, \texttt{mace\_mp}. It runs on the
computer's main processor, the CPU, here on five of its cores at once
(five threads), or on its graphics processor, the GPU, a chip of many
small cores that do the same arithmetic on many numbers at once; and it
keeps its numbers in double precision (float64), the 16 significant
digits of \cref{sec:pe-forces}, or in single precision (float32), about
7. The GPU of the machine used here, an Apple M5 Pro, holds no
double-precision numbers, and the
version of the package used, 0.3.16, stores the energy of each atom in
double precision; the chapter's scripts keep that one output in single
precision on the GPU, which leaves the forces untouched, and check the
result against double precision on the CPU.

LAMMPS runs MACE's own pair styles only when it is built with extra
packages for them, which its standard distribution lacks. Its Python
module offers another way in: the command \texttt{fix ext all external
pf/callback 1 1} makes LAMMPS call a Python function at every step with
the positions, and the forces that function returns are added to the
atoms', its energy to the run's. The function builds an ASE \texttt{Atoms}
from LAMMPS's positions and cell and asks MACE for the forces. On a cell
of 56 atoms of LiC$_6$, eight copies of the seven-atom cell of
\cref{sec:ip-lic6}, moved at random by about \SI{0.05}{\angstrom},
LAMMPS's energy equals ASE's to the last of nine digits, and its forces
equal ASE's to \SI{1.4e-13}{\electronvolt\per\angstrom} once one detail is
handled: LAMMPS holds every cell with its first lattice vector along
$x$ and its second in the $xy$ plane, so ASE turns this cell to fit when
it writes LAMMPS's data file, and a comparison of forces must be made
with the cell as LAMMPS holds it.

\subsection*{A gap: dispersion}

The weak attraction between the sheets of graphite comes from the
correlated motion of their electrons, dispersion, which PBE leaves out.
A model trained on PBE inherits the gap. Graphite's cell holds two
sheets, and its height $c$ is twice their spacing. With MACE-MP-0 alone,
graphite's energy is least at $c = \SI{7.818}{\angstrom}$, 16.5\% above
Trucano and Chen's \SI{6.711}{\angstrom}\cite{TrucanoChen1975}, and
stretching the cell to $c = \SI{9}{\angstrom}$, the sheets
\SI{4.5}{\angstrom} apart, costs only \SI{5.2}{\milli\electronvolt} per
atom (\cref{fig:ip-learned}c). Grimme's D3 correction\cite{Grimme2010}
adds an attraction for every pair of atoms, falling as $1/r^6$ and
$1/r^8$, with coefficients from reference calculations for the 94
elements from hydrogen to plutonium; Becke--Johnson
damping\cite{Grimme2011} keeps it finite where atoms come close. With it
the least energy lies at \SI{6.612}{\angstrom}, 1.5\% below the measured
value, with a binding of \SI{41.2}{\milli\electronvolt} per atom against
$c = \SI{9}{\angstrom}$. Relaxed with the correction, graphite's in-plane
constant is \SI{2.4647}{\angstrom}, 0.03\% above Trucano and Chen's
\SI{2.464}{\angstrom}. The correction is added in ASE by summing its
calculator with MACE's, and it costs \SI{12}{\milli\second} for 126
atoms on the CPU, a twentieth of the medium model's call on the GPU. Every run below includes it.

\subsection*{What it costs}

MACE-MP-0 comes in sizes; the medium model takes \SI{1.20}{\second} for
one call on 126 atoms on five CPU threads in double precision,
\SI{0.62}{\second} in single precision, and \SI{0.25}{\second} on the GPU,
and the small model \num{0.40}, \num{0.21} and \SI{0.081}{\second}
(\cref{fig:ip-learned}a). Every setting grows nearly in proportion to
the number of atoms from 28 to 1008, as the neighbour list makes each
atom's work local, though on the GPU the cost per atom rises by 40\% from
252 to 1008 atoms; for the medium model the GPU is 1.9 to 2.7 times
faster than the CPU in single precision, and 3.5 to 4.8 times faster than
in double. Against the \SI{0.08}{\milli\second} that LAMMPS takes for a step
of 256 argon atoms (\cref{sec:ip-lammps}), a learned potential is
thousands of times dearer per atom, and cheap only beside the electronic
structure it replaces (\cref{sec:ip-aimd}).

\subsection*{The time step}

The test of \cref{sec:in-stability} chooses the step. LiC$_6$, 56 atoms,
was started from velocities drawn at \SI{600}{\kelvin} and run for
\SI{2}{\pico\second} at fixed energy with steps of 0.5, 1, 2 and
\SI{3}{\femto\second} in double precision (\cref{fig:ip-learned}b). The
spread of the total energy is 0.0017, 0.0067, 0.028 and 0.082 of the
kinetic energy's, growing between the shorter steps by 4.02 and 4.16,
against the 4 of $\dt^2$, and from 2 to \SI{3}{\femto\second} by 2.92,
against 2.25: at the longest step the growth outruns $\dt^2$. The mean
total energy of the last quarter of each run differs from that of the
first by less than \SI{0.1}{\milli\electronvolt} up to
\SI{2}{\femto\second}, and by \SI{17.6}{\milli\electronvolt} at
\SI{3}{\femto\second}. Single precision on the GPU gives the same spreads
and the same small differences at 1 and \SI{2}{\femto\second}, so for this
system its rounding is harmless. The longest step that passes is
\SI{2}{\femto\second}; the runs below take \SI{1}{\femto\second}, a
margin for runs of tens of picoseconds at up to \SI{600}{\kelvin}.

\begin{figure}[htb]
\centering
\includegraphics{ch17_practice/learned}
\caption{MACE-MP-0 with D3. (a) The time of one force call against the
number of atoms, LiC$_6$ cells of 28 to 1008 atoms, for the medium model
(solid) and the small (dashed) on five CPU threads in double and single
precision and on the GPU in single precision, and for D3 alone (grey). (b)
LiC$_6$, 56 atoms, at fixed energy: the spread of the total energy over
that of the kinetic energy against the step, in double (circles) and
single (squares) precision near \SI{245}{\kelvin}, and in double
precision near \SI{600}{\kelvin} (triangles), with lines $\propto\dt^2$
through \SI{1}{\femto\second} (dotted). (c) Graphite's energy against the length
$c$ of its cell, with MACE-MP-0 alone and with D3, against the measured
$c$ (dotted).}
\label{fig:ip-learned}
\end{figure}

\begin{keyidea}
A foundation model runs as an ASE calculator, and in LAMMPS through a
Python callback, with the same forces. It inherits the gaps of its
training functional, here dispersion, which D3 restores. It costs about
a quarter of a second per call for a hundred atoms on a GPU, growing
nearly in proportion to the number of atoms, and its time step is chosen by the
same test at fixed energy as any other model's.
\end{keyidea}

\notebook{17}{switch D3 on and off and watch graphite's sheets}

\begin{selfcheck}
A learned potential trained only on bulk metals is used for a metal
surface covered by water. Which part of the system needs checking first,
and how?
\end{selfcheck}
